% Part: normal-modal-logic
% Chapter: filtrations
% Section: preliminaries

\documentclass[../../../include/open-logic-section]{subfiles}

\begin{document}

\olfileid{nml}{fil}{pre}

\olsection{مقدمات}

تصفیه‌ها به ما امکان می‌دهند با نشان‌دادن اینکه دستگاه‌های منطق
موجهاتی‌مان دارای \emph{خاصیت مدل متناهی} هستند، تصمیم‌پذیری آن‌ها را
اثبات کنیم؛ یعنی هر !!{formula} که در مدلی صادق (کاذب) است، در یک مدل
\emph{متناهی} نیز صادق (کاذب) باشد. تصفیه‌ها نسبت به مجموعه‌هایی از
!!{formula}ها تعریف می‌شوند که تحت زیرفرمول‌ها بسته‌اند.

\begin{defn}\ollabel{defn:modallyclosed}
  مجموعهٔ~$\Gamma$ از !!{formula}ها \emph{تحت زیرفرمول‌ها بسته} است
  اگر هر زیرفرمولِ هر !!a{formula} در~$\Gamma$ را شامل شود. افزون‌براین،
  $\Gamma$ \emph{به‌طور موجهاتی بسته} است اگر تحت زیرفرمول‌ها بسته باشد
  و نیز~$!A \in \Gamma$ مستلزم~$\Box!A, \Diamond!A \in \Gamma$ باشد.
\end{defn}

برای نمونه، با داشتن !!a{formula}~$!A$، مجموعهٔ همهٔ
زیر-!!{formula}های آن تحت زیر-!!{formula}ها بسته است. وقتی تصفیهٔ مدلی
را از خلال مجموعهٔ زیر-!!{formula}های~$!A$ تعریف می‌کنیم، خاصیت مطلوب
ما را خواهد داشت: $!A$ را صادق (کاذب) می‌کند اگر و تنها اگر مدل اصلی
چنین کند.

مجموعهٔ جهان‌های تصفیه‌ای از~$\mModel{M}$ از خلال~$\Gamma$ به‌صورت
مجموعهٔ همهٔ رده‌های هم‌ارزیِ رابطهٔ هم‌ارزی زیر تعریف می‌شود.

\begin{defn}
بگذارید~$\mModel{M} =\tuple{W, R, V}$ و فرض کنید~$\Gamma$ تحت
زیر-!!{formula}ها بسته باشد. رابطهٔ~$\equiv$ را روی~$W$ چنان تعریف
کنید که میان هر دو جهانی برقرار باشد که همان !!{formula}های~$\Gamma$
را صادق می‌کنند؛ یعنی:
\[
u \equiv v \quad \text{اگر و تنها اگر} \quad
\forall !A \in \Gamma : \mSat{M}{!A}[u] \Leftrightarrow \mSat{M}{!A}[v].
\]
ردهٔ هم‌ارزیِ~$[w]_\equiv$ برای جهان~$w$، یا به‌اختصار بلوک~$[w]$،
مجموعهٔ همهٔ جهان‌های~$\equiv$-هم‌ارز با~$w$ است:
\[
[w] = \Setabs{v}{v \equiv w}.
\]
\end{defn}

\begin{prop}
  با داشتن~$\mModel{M}$ و~$\Gamma$، رابطهٔ~$\equiv$ چنان‌که در بالا
  تعریف شد، رابطه‌ای هم‌ارزی است؛ یعنی بازتابی، متقارن و متعدی است.
\end{prop}

\begin{proof}
  رابطهٔ~$\equiv$ بازتابی است، زیرا~$w$ دقیقاً همان !!{formula}های
  $\Gamma$ را صادق می‌کند که خودش. متقارن است، زیرا اگر~$u$ همان
  !!{formula}های~$\Gamma$ را صادق کند که~$v$، عکس آن نیز دربارهٔ~$v$
  و~$u$ برقرار است. همچنین متعدی است، زیرا اگر~$u$ همان
  !!{formula}های~$\Gamma$ را صادق کند که~$v$، و~$v$ همان‌ها را که~$w$،
  آنگاه~$u$ همان !!{formula}های~$\Gamma$ را صادق می‌کند که~$w$.
\end{proof}

رابطهٔ~$\equiv$، مانند هر رابطهٔ هم‌ارزی، $W$ را به \emph{رده‌های
هم‌ارزی} یا بلوک‌هایی تقسیم می‌کند؛ یعنی زیرمجموعه‌هایی از~$W$ که
دوبه‌دو جدا هستند و روی‌هم همهٔ~$W$ را می‌پوشانند. هر~$w \in W$،
!!a{element} از یکی از این بلوک‌ها، یعنی~$[w]$، است، زیرا~$w \equiv
w$. پس بلوک‌های~$[w]$ همهٔ~$W$ را می‌پوشانند. آن‌ها دوبه‌دو جدا
هستند، زیرا اگر~$u \in [w]$ و~$u \in [v]$، آنگاه~$u \equiv w$ و
$u \equiv v$؛ و بنا بر تقارن و تعدی~$w \equiv v$، پس~$[w] = [v]$.

\end{document}
